\chapter{Ricci Curvature Comparison}
\label{chap:pet-ricci-curvature-comparison}

\newcommand{\Ric}{\operatorname{Ric}}
\newcommand{\Hess}{\operatorname{Hess}}
\newcommand{\vol}{\operatorname{vol}}
\newcommand{\diam}{\operatorname{diam}}
\newcommand{\Iso}{\mathrm{Iso}}
\newcommand{\snk}{\mathrm{sn}}

Lower Ricci bounds yield volume and Laplacian comparison, analytic
inequalities, diameter rigidity, and splitting results.

\section{Volume Comparison}

\subsection{The Fundamental Equations}

Throughout this section, $(M,g)$ is a complete Riemannian $n$-manifold and
$r(x)=|xp|$ is smooth on $U\subset M$. Write $\partial_r=\nabla r$. The radial
metric equations are
\[
  L_{\partial_r}g=2\Hess r,
  \qquad
  (\nabla_{\partial_r}\Hess r)(X,Y)+\Hess^2r(X,Y)=-R(X,\partial_r,\partial_r,Y),
\]
where $\Hess r(X,Y)=g(\nabla_X\partial_r,Y)$ and
$\Hess^2r(X,Y)=g(\nabla_X\partial_r,\nabla_Y\partial_r)$.

\begin{proposition}[volume-form and Laplacian trace identities]
  \label{prop:pet-ch7-volume-form-laplacian-trace}
  \lean{PetersenLib.volumeFormLaplacianTraceIdentities}
  \source{petersen:page-0291,petersen:page-0292}
  \dcref{ch7:7.1.1}
  The volume form $\vol$ and Laplacian $\Delta r$ of a smooth distance
  function $r$ satisfy
  \[
    \text{(tr1)}\quad L_{\partial_r}\vol=\Delta r\,\vol,
    \qquad
    \text{(tr2)}\quad \partial_r\Delta r+\frac{(\Delta r)^2}{n-1}
    \le\partial_r\Delta r+|\Hess r|^2=-\Ric(\partial_r,\partial_r).
  \]
\end{proposition}
\begin{proof}
  \uses{prop:pet-ch7-volume-form-laplacian-trace}
  \sketch
  The Lie derivative of \(\vol\) along \(\partial_r\) is \(\Delta r\,\vol\). For (tr2), let \(E_1,\ldots,E_n\) be parallel along a radial geodesic. Tracing the radial metric equation gives \(\partial_r\Delta r+|\Hess r|^2=-\Ric(\partial_r,\partial_r)\), while \(\Hess r(\partial_r,\cdot)=0\). Cauchy--Schwarz on the orthogonal \((n-1)\)-space yields \(|\Hess r|^2\ge(\Delta r)^2/(n-1)\).
\end{proof}

Work in exponential polar coordinates at $p$: write $g=dr^2+g_r$ and let
$\vol_{n-1}$ denote the canonical volume form on $S^{n-1}(1)$. The volume form
of $g$ then factors as $\vol=\lambda(r,\theta)\,dr\wedge\vol_{n-1}$,
$\theta\in S^{n-1}$, for a positive function $\lambda$, the \emph{polar
density}. Since $L_{\partial_r}dr=0$ and $L_{\partial_r}\vol_{n-1}=0$ (both
are pulled back from the unit-sphere factor, independent of $r$), applying
(tr1) of \cref{prop:pet-ch7-volume-form-laplacian-trace} gives
$L_{\partial_r}\vol=\partial_r(\lambda)\,dr
\wedge\vol_{n-1}$, so (tr1) reduces to the scalar ODE
$\partial_r\lambda=\lambda\,\Delta r$.

\begin{definition}[polar density and its constant-curvature model]
  \label{def:pet-ch7-polar-density}
  \lean{PetersenLib.polarDensity,PetersenLib.polarDensityConstantCurvature}
  \source{petersen:page-0292,petersen:page-0293}
  \uses{prop:pet-ch7-volume-form-laplacian-trace}
  In exponential polar coordinates $g=dr^2+g_r$ at $p$, the \emph{polar
  density} $\lambda(r,\theta)>0$ is defined by
  $\vol=\lambda\,dr\wedge\vol_{n-1}$; it satisfies
  $\partial_r\lambda=\lambda\,\Delta r$. In the constant curvature $k$ space
  form, $g_k=dr^2+\snk_k^2(r)\,ds_{n-1}^2$, and the corresponding density
  $\lambda_k(r)=\snk_k^{n-1}(r)$ satisfies the same relation with
  $\Delta r=(n-1)\snk_k'(r)/\snk_k(r)$:
  \[
    \partial_r\big(\snk_k^{n-1}(r)\big)
    =(n-1)\frac{\snk_k'(r)}{\snk_k(r)}\,\snk_k^{n-1}(r).
  \]
\end{definition}

\subsection{Volume Estimation}

The trace identities and polar density convert a lower Ricci bound into
radial volume estimates.

\begin{lemma}[Ricci mean-curvature comparison]
  \label{lem:pet-ch7-ricci-comparison-mean-curvature}
  \lean{PetersenLib.ricciMeanCurvatureComparison}
  \source{petersen:page-0293}
  \uses{prop:pet-ch7-volume-form-laplacian-trace, def:pet-ch7-polar-density}
  \dcref{ch7:7.1.2}
  If $(M,g)$ has $\Ric\ge(n-1)k$ for some $k\in\mathbb{R}$, then
  \[
    \Delta r\le(n-1)\frac{\snk_k'(r)}{\snk_k(r)},
    \qquad
    \partial_r\Big(\frac{\lambda}{\lambda_k}\Big)\le0,
    \qquad
    \lambda(r,\theta)\le\lambda_k(r)=\snk_k^{n-1}(r).
  \]
\end{lemma}
\begin{proof}
  \uses{lem:pet-ch7-ricci-comparison-mean-curvature}
  \sketch
  Set \(\rho=\Delta r/(n-1)\). The trace inequality and \(\Ric\ge(n-1)k\) give \(\rho'+\rho^2\le-k\). Its \(1/r+O(r)\) initial behavior agrees with \(\snk_k'/\snk_k\), so Riccati comparison gives the bound for \(\Delta r\). The density ODE in \cref{def:pet-ch7-polar-density} makes \(\lambda/\lambda_k\) nonincreasing, and its limit \(1\) at the origin gives \(\lambda\le\lambda_k\).
\end{proof}

\begin{lemma}[absolute volume comparison]
  \label{lem:pet-ch7-absolute-volume-comparison}
  \lean{PetersenLib.absoluteVolumeComparison}
  \source{petersen:page-0294}
  \uses{lem:pet-ch7-ricci-comparison-mean-curvature}
  \dcref{ch7:7.1.3}
  If $(M,g)$ has $\Ric\ge(n-1)k$, then $\vol B(p,r)\le v(n,k,r)$, where
  $v(n,k,r)$ is the volume of a metric ball of radius $r$ in the constant
  curvature space form $S^n_k$.
\end{lemma}
\begin{proof}
  \uses{lem:pet-ch7-absolute-volume-comparison}
  \sketch
  In polar coordinates, integrate \(\lambda\) over the segment domain of \(B(p,r)\). The pointwise bound \(\lambda\le\lambda_k\) and extension of the nonnegative model density to the whole model ball give \(\vol B(p,r)\le v(n,k,r)\).
\end{proof}

The absolute estimate in \cref{lem:pet-ch7-absolute-volume-comparison}
strengthens to a monotonicity statement for ratios of volumes.

\begin{lemma}[Bishop--Gromov relative volume comparison]
  \label{lem:pet-ch7-relative-volume-comparison}
  \lean{PetersenLib.relativeVolumeComparison}
  \source{petersen:page-0294,petersen:page-0295}
  \uses{lem:pet-ch7-ricci-comparison-mean-curvature}
  \dcref{ch7:7.1.4}
  Let $(M,g)$ be complete with $\Ric\ge(n-1)k$. The volume ratio
  \[
    r\mapsto\frac{\vol B(p,r)}{v(n,k,r)}
  \]
  is a nonincreasing function of $r$, with limit $1$ as $r\to0$.
\end{lemma}
\begin{proof}
  \uses{lem:pet-ch7-relative-volume-comparison}
  \sketch
  By \cref{lem:pet-ch7-ricci-comparison-mean-curvature}, \(N'(R)/D'(R)\) is the spherical average of the nonincreasing ratios \(\lambda/\lambda_k\), where \(N(R)=\vol B(p,R)\) and \(D(R)=v(n,k,R)\). Differentiating \(N/D\) proves monotonicity, and the common polar-density asymptotic gives the limit \(1\) at \(0\).
\end{proof}

\subsection{The Maximum Principle}

Support functions encode one-sided second-derivative bounds for continuous
functions.

\begin{lemma}[Hessian comparison at a contact point]
  \label{lem:pet-ch7-hessian-touching-comparison}
  \lean{PetersenLib.hessianTouchingComparison}
  \source{petersen:page-0295,petersen:page-0296}
  \dcref{ch7:7.1.5}
  If $f,h:(M,g)\to\mathbb{R}$ are $C^2$ functions such that $f(p)=h(p)$ and
  $f(x)\ge h(x)$ for all $x$ near $p$, then
  \[
    \nabla f(p)=\nabla h(p),\qquad \Hess f|_p\ge \Hess h|_p,\qquad \Delta f(p)\ge \Delta h(p).
  \]
\end{lemma}
\begin{proof}
  \uses{lem:pet-ch7-hessian-touching-comparison}
  \sketch
  Apply one-variable calculus to \(f\circ c\) and \(h\circ c\) for geodesics \(c\) through \(p\). Equality of values gives equality of differentials, and the second derivative inequality for every initial velocity gives \(\Hess f|_p\ge\Hess h|_p\); tracing proves the Laplacian inequality.
\end{proof}

\begin{definition}[support functions; barrier-sense Hessian and Laplacian bounds]
  \label{def:pet-ch7-support-function}
  \lean{PetersenLib.SupportFunctionBelow,PetersenLib.barrierHessBound}
  \source{petersen:page-0296}
  \uses{lem:pet-ch7-hessian-touching-comparison}
  By \cref{lem:pet-ch7-hessian-touching-comparison}, a $C^2$ function
  $f:M\to\mathbb{R}$ has $\Hess f|_p\ge B$ for a symmetric bilinear form $B$ on $T_pM$
  (respectively $\Delta f(p)\ge a\in\mathbb{R}$) if and only if for every
  $\varepsilon>0$ there is a function $f_\varepsilon$ defined near $p$, called a
  \emph{support function from below} for $f$ at $p$, with
  \begin{enumerate}
    \item[(1)] $f_\varepsilon(p)=f(p)$;
    \item[(2)] $f(x)\ge f_\varepsilon(x)$ near $p$;
    \item[(3)] $\Hess f_\varepsilon|_p\ge B-\varepsilon g|_p$ (respectively
      $\Delta f_\varepsilon(p)\ge a-\varepsilon$).
  \end{enumerate}
  \emph{Support functions from above} are defined symmetrically and give upper
  bounds for $\Hess f$ and $\Delta f$. For a merely continuous $f:(M,g)\to\mathbb{R}$
  one \emph{defines} $\Hess f|_p\ge B$ (respectively $\Delta f(p)\ge a$), said to
  hold in the \emph{support} or \emph{barrier sense}, to mean that such smooth
  support functions $f_\varepsilon$ from below exist for every $\varepsilon>0$.
  A continuous $f:(M,g)\to\mathbb{R}$ is \emph{convex} if $\Hess f\ge0$ everywhere in
  this sense; this matches the classical definition when $(M,g)\subset(\mathbb{R},g_\mathbb{R})$.
\end{definition}

\begin{theorem}[convex functions near a maximum]
  \label{thm:pet-ch7-convex-constant-near-max}
  \lean{PetersenLib.convexHessNonnegConstantNearMax}
  \source{petersen:page-0296}
  \uses{def:pet-ch7-support-function}
  \dcref{ch7:7.1.6}
  If $f:(M,g)\to\mathbb{R}$ is continuous with $\Hess f\ge0$ everywhere, then $f$ is
  constant near any local maximum. In particular, $f$ cannot have a global
  maximum unless $f$ is constant.
\end{theorem}

\begin{definition}[subharmonic and superharmonic functions]
  \label{def:pet-ch7-subharmonic-superharmonic}
  \lean{PetersenLib.IsSubharmonic,PetersenLib.IsSuperharmonic}
  \source{petersen:page-0296}
  \uses{def:pet-ch7-support-function}
  A continuous $f:(M,g)\to\mathbb{R}$ with $\Delta f\ge0$ everywhere (in the support
  sense) is \emph{subharmonic}; if $\Delta f\le0$ it is \emph{superharmonic}.
\end{definition}

\begin{theorem}[strong maximum principle]
  \label{thm:pet-ch7-strong-maximum-principle}
  \lean{PetersenLib.strongMaximumPrinciple}
  \source{petersen:page-0296,petersen:page-0297,petersen:page-0298}
  \uses{def:pet-ch7-subharmonic-superharmonic, thm:pet-ch7-convex-constant-near-max}
  \dcref{ch7:7.1.7}
  If $f:(M,g)\to\mathbb{R}$ is continuous and subharmonic, then $f$ is constant in a
  neighborhood of every local maximum. In particular, if $f$ has a global
  maximum then $f$ is constant.
\end{theorem}
\begin{proof}
  \uses{thm:pet-ch7-strong-maximum-principle,def:pet-ch7-support-function,lem:pet-ch7-hessian-touching-comparison}
  \sketch
  If \(\Delta f>0\), a support function from below at a local maximum has nonpositive Hessian trace, a contradiction. For \(\Delta f\ge0\), perturb \(f\) on a small ball by \(\delta(e^{\alpha\phi}-1)\) with positive Laplacian and boundary values below the maximum. An interior maximum of the perturbed function produces the same contradiction via a support function. Thus \(f\) is locally constant at every maximum.
\end{proof}

\begin{definition}[linear and harmonic functions]
  \label{def:pet-ch7-linear-harmonic-function}
  \lean{PetersenLib.IsLinearFunction,PetersenLib.IsHarmonicFunction}
  \source{petersen:page-0298}
  \uses{def:pet-ch7-support-function}
  A continuous $f:(M,g)\to\mathbb{R}$ is \emph{linear} if $\Hess f\equiv0$, i.e.\
  both $\Hess f\ge0$ and $\Hess f\le0$ hold everywhere in the support sense.
  Then $f\circ c$ is both convex and concave along every geodesic $c$, so
  $(f\circ c)(t)=f(c(0))+at$; equivalently
  $f\circ\exp_p(x)=f(p)+g(v_p,x)$ for a fixed $v_p\in T_pM$, and $f$ is
  automatically $C^\infty$ with $\nabla f|_p=v_p$. More generally, $f$ is
  \emph{harmonic} if it is continuous with $\Delta f=0$ (in the support sense).
\end{definition}

\begin{theorem}[regularity of weakly harmonic functions]
  \label{thm:pet-ch7-harmonic-regularity}
  \lean{PetersenLib.harmonicFunctionRegularity}
  \source{petersen:page-0298,petersen:page-0299}
  \uses{def:pet-ch7-linear-harmonic-function}
  \dcref{ch7:7.1.8}
  If $f:(M,g)\to\mathbb{R}$ is continuous and harmonic in the weak (support) sense,
  then $f$ is smooth.
\end{theorem}
\begin{proof}
  \uses{thm:pet-ch7-harmonic-regularity,thm:pet-ch7-strong-maximum-principle}
  \sketch
  Solve the smooth Dirichlet problem on a coordinate domain with boundary data \(f\). Both \(u-f\) and \(f-u\) are subharmonic and vanish on the boundary. The strong maximum principle, \cref{thm:pet-ch7-strong-maximum-principle}, forces both to vanish, so \(f=u\) and is smooth.
\end{proof}

\subsection{Geometric Laplacian Comparison}

The triangle inequality supplies upper support functions for distance
functions.

\begin{lemma}[Calabi's Laplacian comparison]
  \label{lem:pet-ch7-laplacian-comparison-calabi}
  \lean{PetersenLib.laplacianComparisonCalabi}
  \source{petersen:page-0298,petersen:page-0299}
  \uses{def:pet-ch7-support-function}
  \dcref{ch7:7.1.9}
  If $(M,g)$ is complete and $\Ric(M,g)\ge(n-1)k$, then any distance function
  $r(x)=|xp|$ satisfies
  \[
    \Delta r(x)\le(n-1)\frac{\snk_k'(r(x))}{\snk_k(r(x))}.
  \]
\end{lemma}
\begin{proof}
  \uses{lem:pet-ch7-laplacian-comparison-calabi,lem:pet-ch7-ricci-comparison-mean-curvature}
  \sketch
  At \(q\), use \(r_\varepsilon(x)=\varepsilon+|\sigma(\varepsilon)x|\) as an upper support for \(r=|px|\). Apply the smooth comparison theorem to \(r_\varepsilon\), then let \(\varepsilon\to0\); the possible cut or critical-point alternatives are excluded by the minimizing segment from \(p\) to \(q\).
\end{proof}

\subsection{The Segment, Poincaré, and Sobolev Inequalities}

For a bounded domain $B\subset M$ and $u:M\to\mathbb{R}$, write
\[
  \|u\|_{p,B}=\left(\frac{1}{\vol B}\int_B|u|^p\,\vol\right)^{1/p},\qquad
  u_B=\frac{1}{\vol B}\int_Bu\,\vol.
\]

\begin{theorem}[Cheeger--Colding segment inequality]
  \label{thm:pet-ch7-segment-inequality}
  \lean{PetersenLib.segmentInequality}
  \source{petersen:page-0300}
  \dcref{ch7:7.1.10}
  Assume $(M,g)$ has $\Ric\ge(n-1)k$, $k\le0$. Let $f:M\to[0,\infty)$ and
  $A,B\subset W\subset M$ with $\diam W\le D$. For each $x\in A$, $y\in B$ fix a
  segment $c_{x,y}:[0,1]\to M$ from $x$ to $y$, and suppose $c_{x,y}(t)\in W$
  for all $x\in A$, $y\in B$, $t\in[0,1]$. Then
  \[
    \int_A\int_B\int_0^1 f\circ c_{x,y}(t)\,dt\,\vol_x\,\vol_y
    \le C(\vol A+\vol B)\int_Wf\,\vol,
  \]
  where $C=C(n,kD^2)$.
\end{theorem}
\begin{proof}
  \uses{thm:pet-ch7-segment-inequality,lem:pet-ch7-ricci-comparison-mean-curvature}
  \sketch
  For \(t\ge\frac12\), polar coordinates at \(x\) and density comparison bound the Jacobian ratio along \(y\mapsto c_{x,y}(t)\) by \(C(n,kD^2)\), yielding a \(C\vol A\) estimate. The symmetric estimate for \(t\le\frac12\) contributes \(C\vol B\); adding them proves the segment inequality.
\end{proof}

\begin{corollary}[weak Poincaré inequality]
  \label{cor:pet-ch7-weak-poincare}
  \lean{PetersenLib.weakPoincareInequality}
  \source{petersen:page-0301,petersen:page-0302}
  \uses{thm:pet-ch7-segment-inequality}
  \dcref{ch7:7.1.11}
  Assume $(M,g)$ has $\Ric\ge(n-1)k$, $k\le0$. Any smooth $u:M\to[0,\infty)$
  satisfies
  \[
    \|u-u_{B(p,R)}\|_{1,B(p,R)}\le4C^2R\,\|du\|_{1,B(p,2R)},
  \]
  where $R\le D$ and $C=C(n,kD^2)$ is as in \cref{thm:pet-ch7-segment-inequality}.
\end{corollary}
\begin{proof}
  \uses{cor:pet-ch7-weak-poincare,thm:pet-ch7-segment-inequality,lem:pet-ch7-relative-volume-comparison}
  \sketch
  For \(B=B(p,R)\), integrate \(|u(x)-u(y)|\) along chosen minimizing segments in \(B(p,2R)\). The segment inequality bounds the triple integral of \(|du|\), and \cref{lem:pet-ch7-relative-volume-comparison} bounds \(\vol B(p,2R)/\vol B(p,R)\), giving the stated \(4C^2R\) estimate.
\end{proof}

\begin{remark}[upper-gradient form of weak Poincaré]
  \label{rem:pet-ch7-upper-gradient}
  \lean{PetersenLib.upperGradientPoincare}
  \source{petersen:page-0302}
  \uses{cor:pet-ch7-weak-poincare}
  \dcref{ch7:7.1.12}
  \Cref{cor:pet-ch7-weak-poincare} holds for any measurable $u$ with a function
  $G$ in place of $|du|$, provided $|u(x)-u(y)|\le\int_0^1G(c(t))\,|\dot
  c(t)|\,dt$ for all $c\in\Omega_{x,y}$ (paths from $x$ to $y$); such $G$ is an
  \emph{upper gradient} for $u$.
\end{remark}

\begin{definition}[maximal function]
  \label{def:pet-ch7-maximal-function}
  \lean{PetersenLib.maximalFunction}
  \source{petersen:page-0302}
  For $u:M\to\mathbb{R}$ and $D>0$, the \emph{maximal function} of $u$ (relative to
  the scale $D$) is
  \[
    M(u)(x)=\sup_{R\in(0,D]}\frac{1}{\vol B(x,R)}\int_{B(x,R)}|u|\,\vol.
  \]
\end{definition}

\begin{theorem}[maximal-function weak estimate]
  \label{thm:pet-ch7-maximal-function-theorem}
  \lean{PetersenLib.maximalFunctionTheorem}
  \source{petersen:page-0303}
  \uses{def:pet-ch7-maximal-function}
  \dcref{ch7:7.1.14}
  Assume $(M,g)$ has $\Ric\ge(n-1)k$. There is $C=C(n,kD^2)$ such that for all
  $t>0$,
  \[
    t\,\vol\{M(u)>t\}\le C\int|u|\,\vol.
  \]
\end{theorem}
\begin{proof}
  \uses{thm:pet-ch7-maximal-function-theorem,lem:pet-ch7-relative-volume-comparison,rem:pet-ch7-ex-5}
  \sketch
  For each \(x\) with \(M(u)(x)>t\), choose \(R_x\) with \(t\vol B(x,R_x)<\int_{B(x,R_x)}|u|\). The covering lemma \cref{rem:pet-ch7-ex-5} supplies disjoint \(B(x_i,R_i)\) whose fivefold enlargements cover the level set. Relative volume comparison and disjointness give \(t\vol\{M(u)>t\}\le C\int|u|\).
\end{proof}

\begin{theorem}[weak Poincaré--Sobolev level-set estimate]
  \label{thm:pet-ch7-weak-poincare-sobolev-estimate}
  \lean{PetersenLib.weakPoincareSobolevEstimate}
  \source{petersen:page-0303,petersen:page-0304,petersen:page-0305}
  \uses{cor:pet-ch7-weak-poincare, thm:pet-ch7-maximal-function-theorem}
  \dcref{ch7:7.1.15}
  Assume $(M,g)$ has $\Ric\ge(n-1)k$, $k\le0$, and $\diam M\le D$. For smooth
  $u:M\to[0,\infty)$,
  \[
    t^{\frac{n}{n-1}}\,\vol\{u-u_{B(x,R)}>t\}
    \le C R^{\frac{n}{n-1}}\,\vol M\,\|du\|_1^{\frac{n}{n-1}},
  \]
  where $C=C(n,kD^2)$.
\end{theorem}
\begin{proof}
  \uses{thm:pet-ch7-weak-poincare-sobolev-estimate,cor:pet-ch7-weak-poincare,thm:pet-ch7-maximal-function-theorem,lem:pet-ch7-relative-volume-comparison}
  \sketch
  Use dyadic balls \(B_i=B(x,2^{-i}D)\). Telescoping averages and weak Poincare reduce pointwise oscillation to \(\sum_iR_i\|du\|_{1,B_{i-1}}\). Split at a scale \(r\), bounding the tail by \(M(|du|)\) and the head by relative volume comparison. Balancing \(r\) and applying \cref{thm:pet-ch7-maximal-function-theorem} yields the asserted level-set bound.
\end{proof}

\begin{theorem}[Poincaré--Sobolev inequality]
  \label{thm:pet-ch7-poincare-sobolev}
  \lean{PetersenLib.poincareSobolevInequality}
  \source{petersen:page-0302,petersen:page-0305,petersen:page-0306,petersen:page-0307}
  \uses{thm:pet-ch7-weak-poincare-sobolev-estimate}
  \dcref{ch7:7.1.13}
  Assume $(M,g)$ has $\Ric\ge(n-1)k$, $k\le0$. For all smooth $u:M\to[0,\infty)$
  and $\nu\in[1,\tfrac{n}{n-1}]$,
  \[
    \|u-u_{B(x,R)}\|_{\nu,B(x,R)}\le C(n,kD^2)\,R\,\|du\|_{1,B(x,R)},
  \]
  where $R\le D$.
\end{theorem}
\begin{proof}
  \uses{thm:pet-ch7-poincare-sobolev,thm:pet-ch7-weak-poincare-sobolev-estimate}
  \sketch
  Reduce to a nonnegative function vanishing on at least half of \(M\) by subtracting a median and splitting into positive and negative parts. Apply \cref{thm:pet-ch7-weak-poincare-sobolev-estimate} to dyadic truncations, sum over disjoint level bands, and use \(\ell^{n/(n-1)}\)-summability of their gradient masses. This gives the \(n/(n-1)\) estimate; Hölder gives the smaller exponents.
\end{proof}

\begin{remark}[Dirichlet Poincaré inequality]
  \label{rem:pet-ch7-dirichlet-poincare}
  \lean{PetersenLib.exerciseDirichletPoincareReference}
  \source{petersen:page-0307}
  \uses{thm:pet-ch7-poincare-sobolev}
  \dcref{ch7:7.1.16}
  See \cref{rem:pet-ch7-ex-17} for the Poincaré inequality for functions with
  Dirichlet boundary conditions vanishing on a ball's boundary.
\end{remark}

\begin{proposition}[Sobolev interpolation hierarchy]
  \label{prop:pet-ch7-sobolev-hierarchy}
  \lean{PetersenLib.sobolevInterpolationHierarchy}
  \source{petersen:page-0307}
  \uses{thm:pet-ch7-poincare-sobolev}
  \dcref{ch7:7.1.17}
  Suppose every smooth function on $(M,g)$ satisfies $\|u-u_M\|_{s-1}\le
  S\|du\|_1$ for some fixed $s>1$. Then for $1\le p<s$,
  \[
    \|u\|_{\frac{sp}{s-p}}\le\frac{p(s-1)}{s-p}S\|du\|_p+\|u\|_p.
  \]
\end{proposition}
\begin{proof}
  \uses{prop:pet-ch7-sobolev-hierarchy}
  \sketch
  For \(p=1\), compare \(u\) with its mean and use \(|u_M|\le\|u\|_1\). For \(p>1\), apply the assumed estimate to \(|u|^q\), use Hölder on \(q|u|^{q-1}|du|\), and choose \(q=p(s-1)/(s-p)\). Dividing by the resulting norm gives the claimed coefficient.
\end{proof}

\begin{theorem}[Rellich compactness]
  \label{thm:pet-ch7-rellich-compactness}
  \lean{PetersenLib.rellichCompactness}
  \source{petersen:page-0308}
  \uses{thm:pet-ch7-weak-poincare-sobolev-estimate}
  \dcref{ch7:7.1.18}
  Assume $(M^n,g)$ is a compact Riemannian $n$-manifold. Then the inclusion
  $W^{1,2}(M)\subset L^2(M)$ (where $W^{1,2}(M)$ is the Hilbert-space closure
  of $C^\infty(M)$ in the norm $\|u\|_2^2+\|du\|_2^2$) is compact.
\end{theorem}
\begin{proof}
  \uses{thm:pet-ch7-rellich-compactness,thm:pet-ch7-weak-poincare-sobolev-estimate}
  \sketch
  A bounded sequence in \(W^{1,2}\) has a weakly convergent subsequence. The weak Poincare--Sobolev level estimate gives convergence in measure and then almost everywhere along a further subsequence. The bounds obtained from \cref{thm:pet-ch7-poincare-sobolev} and \cref{prop:pet-ch7-sobolev-hierarchy} give uniform integrability, hence convergence in \(L^2\).
\end{proof}

\section{Applications of Ricci Curvature Comparison}

\subsection{Finiteness of Fundamental Groups}

Volume comparison bounds short presentations of fundamental groups.

\begin{lemma}[Gromov's short presentation of the fundamental group]
  \label{lem:pet-ch7-gromov-generators}
  \lean{PetersenLib.gromovShortGenerators}
  \source{petersen:page-0308,petersen:page-0309}
  \dcref{ch7:7.2.1}
  A compact Riemannian manifold $M$ admits generators $\{c_1,\dots,c_m\}$ for
  $\Gamma=\pi_1(M)$ such that all relations are of the form
  $c_ic_jc_k^{-1}=1$ for suitable $i,j,k$, with each generator represented by a
  loop of length $\le3\diam(M)$.
\end{lemma}
\begin{proof}
  \uses{lem:pet-ch7-gromov-generators}
  \sketch
  Choose a sufficiently fine triangulation. Loops formed by a basepoint path, one edge, and the reverse basepoint path generate the fundamental group; boundaries of 2-simplices give all relations \(c_ic_jc_k^{-1}=1\). Their lengths are at most \(2\diam M+\varepsilon\le3\diam M\).
\end{proof}

\begin{example}[lens spaces with collapsing volume]
  \label{ex:pet-ch7-lens-spaces-example}
  \lean{PetersenLib.exampleLensSpacesGrowingFundamentalGroup}
  \source{petersen:page-0309}
  \uses{lem:pet-ch7-gromov-generators}
  \dcref{ch7:7.2.2}
  Consider $M_k=S^3/\mathbb{Z}_k$, the constant curvature $3$-sphere quotient by
  the cyclic group of order $k$. As $k\to\infty$, $\vol M_k\to0$ while the
  curvature stays $1$ and $\diam M_k=\pi/2$. Thus $\pi_1(M_k)=\mathbb{Z}_k$ grows
  only at the expense of shrinking volume; writing $\mathbb{Z}_k$ via generators as
  in \cref{lem:pet-ch7-gromov-generators} requires the number of generators to
  go to infinity with $k$, foreshadowing \cref{thm:pet-ch7-anderson-finiteness}.
\end{example}

\begin{definition}[the class $\mathfrak{M}(n,k,v,D)$]
  \label{def:pet-ch7-manifold-class}
  \lean{PetersenLib.ManifoldClass}
  \source{petersen:page-0309}
  For $n\in\mathbb{N}$, $k\in\mathbb{R}$, $v,D\in(0,\infty)$, let
  $\mathfrak{M}(n,k,v,D)$ denote the class of compact Riemannian $n$-manifolds
  $(M,g)$ with $\Ric\ge(n-1)k$, $\vol M\ge v$, $\diam M\le D$.
\end{definition}

\begin{theorem}[Anderson's fundamental-group finiteness theorem]
  \label{thm:pet-ch7-anderson-finiteness}
  \lean{PetersenLib.andersonFundamentalGroupFiniteness}
  \source{petersen:page-0309}
  \uses{def:pet-ch7-manifold-class, lem:pet-ch7-gromov-generators}
  \dcref{ch7:7.2.3}
  For fixed $n,k,v,D$, there are only finitely many isomorphism types of
  fundamental group among the manifolds in $\mathfrak{M}(n,k,v,D)$.
\end{theorem}
\begin{proof}
  \uses{thm:pet-ch7-anderson-finiteness,lem:pet-ch7-gromov-generators,lem:pet-ch7-absolute-volume-comparison}
  \sketch
  Use short generators and a Dirichlet fundamental domain \(F\) in the universal cover. The translates of \(F\) associated to distinct generators are disjoint and lie in \(B(x,6D)\). Absolute volume comparison bounds their number by \(v(n,k,6D)/v\), so only finitely many relation presentations occur.
\end{proof}

\begin{lemma}[Anderson's short-loop subgroup bound]
  \label{lem:pet-ch7-anderson-short-loops}
  \lean{PetersenLib.andersonShortLoopSubgroupFinite}
  \source{petersen:page-0310}
  \uses{def:pet-ch7-manifold-class}
  \dcref{ch7:7.2.4}
  For fixed $n,k,v,D$ there are $L=L(n,k,v,D)$, $N=N(n,k,v,D)$ such that for
  every $M\in\mathfrak{M}(n,k,v,D)$, any subgroup of $\pi_1(M)$ generated by
  loops of length $\le L$ has order $\le N$.
\end{lemma}
\begin{proof}
  \uses{lem:pet-ch7-anderson-short-loops,lem:pet-ch7-absolute-volume-comparison}
  \sketch
  For a subgroup generated by loops of length at most \(L\), products of at most \(r\) generators move \(F\) into \(B(x,rL+D)\). Disjoint translates and absolute comparison bound the number of such products. Choosing \(L\) and \(N\) from the resulting bound forces every generated subgroup to have order at most \(N\).
\end{proof}

\subsection{Maximal Diameter Rigidity}

Equality in Myers' diameter bound is rigid.

\begin{theorem}[Cheng's maximal-diameter rigidity theorem]
  \label{thm:pet-ch7-cheng-rigidity}
  \lean{PetersenLib.chengMaximalDiameterRigidity}
  \source{petersen:page-0310,petersen:page-0311,petersen:page-0312}
  \uses{lem:pet-ch7-laplacian-comparison-calabi, thm:pet-ch7-strong-maximum-principle}
  \dcref{ch7:7.2.5}
  If $(M,g)$ is complete with $\Ric\ge(n-1)k>0$ and $\diam=\pi/\sqrt k$, then
  $(M,g)$ is isometric to $S^n_k$ (the space form of curvature $k$).
\end{theorem}
\begin{proof}
  \uses{thm:pet-ch7-cheng-rigidity,lem:pet-ch7-laplacian-comparison-calabi,thm:pet-ch7-strong-maximum-principle,prop:pet-ch7-volume-form-laplacian-trace}
  \sketch
  For points \(p,q\) at distance \(\pi/\sqrt{k}\), the excess \(r+\widetilde r\) is at least this value and has nonpositive Laplacian by comparison, so the strong maximum principle makes it constant. Concatenating minimizing segments through any other point is smooth, so both distance functions are smooth there. Equality in the Riccati and Cauchy--Schwarz comparisons gives \(\Hess r=(\snk_k'/\snk_k)g_r\); integrating the radial ODE yields \(g=dr^2+\snk_k^2(r)ds_{n-1}^2\) and \(M\cong S_k^n\).
\end{proof}

\begin{example}[almost maximal diameter in dimension four]
  \label{ex:pet-ch7-anderson-cp2-example}
  \lean{PetersenLib.andersonAlmostMaximalDiameterExampleFour}
  \source{petersen:page-0312,petersen:page-0313}
  \uses{thm:pet-ch7-cheng-rigidity}
  \dcref{ch7:7.2.6}
  For $n=4$, consider metrics $dr^2+\rho^2\sigma_1^2+\phi^2(\sigma_2^2+\sigma_3^2)$
  on $I\times S^3$, with
  \[
    \rho(r)=\begin{cases}\dfrac{\sin(ar)}{a}&r\le r_0,\\c_1\sin(r+\delta)&r\ge
    r_0,\end{cases}\qquad
    \phi(r)=\begin{cases}br^2+c&r\le r_0,\\c_2\sin(r+\delta)&r\ge
    r_0,\end{cases}
  \]
  reflected in $r=\pi/2-\delta$ to produce a metric on $\mathbb{CP}^2\#\overline{\mathbb{CP}^2}$.
  For any small $r_0>0$ the parameters can be tuned so $\rho,\phi$ are $C^1$
  and generate a metric with $\Ric\ge3$; as $r_0\to0$ the parameter interval
  $I\to[0,\pi]$, so the diameters converge to $\pi$, showing that
  \cref{thm:pet-ch7-cheng-rigidity} does not admit a topological rigidity
  refinement when $n=4$.
\end{example}

\begin{example}[higher-dimensional almost maximal diameter]
  \label{ex:pet-ch7-otsu-example}
  \lean{PetersenLib.otsuAlmostMaximalDiameterExample}
  \source{petersen:page-0313}
  \uses{ex:pet-ch7-anderson-cp2-example}
  \dcref{ch7:7.2.7}
  For $n\ge5$, standard doubly warped products $dr^2+\rho^2ds_2^2+\phi^2ds^2_{n-3}$
  on $I\times S^2\times S^{n-3}$, with $\rho,\phi$ chosen analogously to
  \cref{ex:pet-ch7-anderson-cp2-example}, yield metrics on $S^2\times
  S^{n-2}$ with $\Ric\ge n-1$ and diameter $\to\pi$. In both examples the
  functions $\rho,\phi$ are only $C^1$ at the outset, but being concave they
  can be smoothed near the break points while remaining concave, changing
  values and first derivatives negligibly and only increasing the (already
  nonpositive) second derivative in absolute value, so the lower Ricci bound
  persists after smoothing.
\end{example}

\section{Manifolds of Nonnegative Ricci Curvature}

Throughout this section, $(M,g)$ is complete and noncompact.

\subsection{Rays and Lines}

\begin{definition}[ray, line]
  \label{def:pet-ch7-ray-line}
  \lean{PetersenLib.IsRay,PetersenLib.IsLine}
  \source{petersen:page-0313}
  A \emph{ray} $r:[0,\infty)\to(M,g)$ is a unit-speed geodesic with
  $|r(t)r(s)|=|t-s|$ for all $t,s\ge0$; it may be viewed as a semi-infinite
  segment, or a segment from $r(0)$ to infinity. A \emph{line}
  $l:\mathbb{R}\to(M,g)$ is a unit-speed geodesic with $|l(t)l(s)|=|t-s|$ for all
  $t,s\in\mathbb{R}$.
\end{definition}

\begin{definition}[connected/disconnected at infinity]
  \label{def:pet-ch7-connected-infinity}
  \lean{PetersenLib.IsConnectedAtInfinity}
  \source{petersen:page-0314}
  A complete manifold $(M,g)$ is \emph{connected at infinity} if for every
  compact $K\subset M$ there is a compact $C\supset K$ such that any two
  points of $M-C$ can be joined by a curve in $M-K$. Otherwise $(M,g)$ is
  \emph{disconnected at infinity}.
\end{definition}

\begin{lemma}[existence of rays and lines]
  \label{lem:pet-ch7-ray-line-existence}
  \lean{PetersenLib.rayExistenceLineIfDisconnectedAtInfinity}
  \source{petersen:page-0313,petersen:page-0314}
  \uses{def:pet-ch7-ray-line, def:pet-ch7-connected-infinity}
  \dcref{ch7:7.3.1}
  If $p\in(M,g)$ then a ray emanating from $p$ always exists. If $M$ is
  disconnected at infinity, $(M,g)$ contains a line.
\end{lemma}
\begin{proof}
  \uses{lem:pet-ch7-ray-line-existence}
  \sketch
  A ray from \(p\) is obtained by taking a limit of minimizing segments to points escaping to infinity. If the manifold is disconnected at infinity, choose segments between escaping points that cross a fixed compact set; after recentering at the crossing points, a subsequential limit is a line.
\end{proof}

\begin{example}[paraboloid metrics without lines]
  \label{ex:pet-ch7-paraboloid-no-lines}
  \lean{PetersenLib.paraboloidContainsNoLines}
  \source{petersen:page-0314}
  \uses{def:pet-ch7-ray-line}
  \dcref{ch7:7.3.2}
  Surfaces of revolution $dr^2+\rho^2(r)ds^2_{n-1}$ with $\rho:[0,\infty)\to
  [0,\infty)$, $\dot\rho(t)<1$, $\ddot\rho(t)<0$ for $t>0$, contain no lines;
  such manifolds look like paraboloids.
\end{example}

\begin{example}[cylindrical manifolds contain a line]
  \label{ex:pet-ch7-cylinder-contains-line}
  \lean{PetersenLib.sphereCrossRAlwaysContainsLine}
  \source{petersen:page-0315}
  \uses{lem:pet-ch7-ray-line-existence}
  \dcref{ch7:7.3.3}
  Any complete metric on $S^{n-1}\times\mathbb{R}$ contains a line, since the
  manifold is disconnected at infinity.
\end{example}

\begin{example}[Schwarzschild metrics contain no line]
  \label{ex:pet-ch7-schwarzschild-no-line}
  \lean{PetersenLib.schwarzschildContainsNoLine}
  \source{petersen:page-0315}
  \uses{def:pet-ch7-ray-line}
  \dcref{ch7:7.3.4}
  The Schwarzschild metric on $S^{n-2}\times\mathbb{R}^2$ contains no lines, as
  follows from the splitting theorem (\cref{thm:pet-ch7-splitting-theorem})
  since the space is not metrically a product.
\end{example}

\subsection{Busemann Functions}

Assume $\Ric\ge0$. For a unit-speed ray $c:[0,\infty)\to(M,g)$, define
$b_t(x)=|xc(t)|-t$.

\begin{proposition}[estimates for Busemann approximants]
  \label{prop:pet-ch7-busemann-approximants}
  \lean{PetersenLib.busemannApproximantsProperties}
  \source{petersen:page-0316}
  \uses{def:pet-ch7-ray-line, lem:pet-ch7-laplacian-comparison-calabi}
  \dcref{ch7:7.3.6}
  The functions $b_t$, $t\ge0$, satisfy:
  \begin{enumerate}
    \item[(1)] for fixed $x$, $t\mapsto b_t(x)$ is decreasing and bounded in
      absolute value by $|xc(0)|$;
    \item[(2)] $|b_t(x)-b_t(y)|\le|xy|$;
    \item[(3)] $\Delta b_t(x)\le\dfrac{n-1}{b_t(x)+t}$ everywhere.
  \end{enumerate}
\end{proposition}
\begin{proof}
  \uses{prop:pet-ch7-busemann-approximants}
  \sketch
  The triangle inequality gives \(|b_t(x)|\le|xc(0)|\), monotonicity in \(t\), and the Lipschitz bound. Since \(b_t+t\) is a distance function based at \(c(t)\), \cref{lem:pet-ch7-laplacian-comparison-calabi} gives \(\Delta b_t\le(n-1)/(b_t+t)\).
\end{proof}

\begin{definition}[Busemann function]
  \label{def:pet-ch7-busemann-function}
  \lean{PetersenLib.busemannFunction}
  \source{petersen:page-0316}
  \uses{prop:pet-ch7-busemann-approximants}
  By \cref{prop:pet-ch7-busemann-approximants}, $b_t$ converges pointwise as
  $t\to\infty$ to a $1$-Lipschitz function
  \[
    b_\infty(x)=\lim_{t\to\infty}b_t(x),\qquad |b_\infty(x)-b_\infty(y)|\le|xy|,
    \qquad |b_\infty(x)|\le|xc(0)|,
  \]
  satisfying $b_\infty(c(r))=\lim_{t\to\infty}(|c(r)c(t)|-t)=-r$. This
  $b_\infty$, denoted $b_c$, is the \emph{Busemann function} for $c$, thought
  of as a renormalized distance function from ``$c(\infty)$.''
\end{definition}

\begin{example}[Euclidean Busemann functions]
  \label{ex:pet-ch7-euclidean-busemann}
  \lean{PetersenLib.euclideanBusemannFunction}
  \source{petersen:page-0316}
  \uses{def:pet-ch7-busemann-function}
  \dcref{ch7:7.3.7}
  If $M=(\mathbb{R}^n,g_{\mathbb{R}^n})$, every Busemann function has the form
  $b_c(x)=\dot c(0)\cdot(c(0)-x)$.
\end{example}

\begin{definition}[horosphere, asymptote]
  \label{def:pet-ch7-horosphere-asymptote}
  \lean{PetersenLib.Horosphere,PetersenLib.IsAsymptoteFor}
  \source{petersen:page-0317}
  \uses{def:pet-ch7-busemann-function}
  The level sets $b_c^{-1}(t)$ are \emph{horospheres} (hyperplanes in $\mathbb{R}^n$;
  spheres tangent to the boundary in the Poincaré model of hyperbolic space).
  Given a ray $c$ and $p\in M$, a family of unit-speed segments
  $\sigma_i:[0,L_i]\to M$ from $p$ to $c(t_i)$, $t_i\to\infty$, subconverges (as
  in \cref{lem:pet-ch7-ray-line-existence}) to a ray $\tilde c:[0,\infty)\to
  M$, $\tilde c(0)=p$, called an \emph{asymptote} for $c$ from $p$; asymptotes
  need not be unique.
\end{definition}

\begin{proposition}[Busemann comparison along an asymptote]
  \label{prop:pet-ch7-busemann-asymptote-relation}
  \lean{PetersenLib.busemannAsymptoteRelation}
  \source{petersen:page-0317,petersen:page-0318}
  \uses{def:pet-ch7-horosphere-asymptote}
  \dcref{ch7:7.3.8}
  With $c,p,\tilde c$ as in \cref{def:pet-ch7-horosphere-asymptote}, the
  Busemann functions satisfy
  \begin{enumerate}
    \item[(1)] $b_c(x)\le b_c(p)+b_{\tilde c}(x)$ for all $x\in M$;
    \item[(2)] $b_c(\tilde c(t))=b_c(p)+b_{\tilde c}(\tilde c(t))=b_c(p)-t$ for
      all $t\ge0$.
  \end{enumerate}
\end{proposition}
\begin{proof}
  \uses{prop:pet-ch7-busemann-asymptote-relation}
  \sketch
  For segments \(\sigma_i\) from \(p\) to \(c(t_i)\) converging to an asymptote \(\widetilde c\), split their lengths and use \cref{def:pet-ch7-busemann-function} to obtain \(b_c(\widetilde c(s))=b_c(p)-s\). Applying the triangle inequality through \(\widetilde c(t)\) and then letting \(t\to\infty\) gives \(b_c(x)\le b_c(p)+b_{\widetilde c}(x)\).
\end{proof}

\begin{lemma}[superharmonicity of Busemann functions]
  \label{lem:pet-ch7-busemann-superharmonic}
  \lean{PetersenLib.busemannFunctionSuperharmonic}
  \source{petersen:page-0318}
  \uses{prop:pet-ch7-busemann-asymptote-relation, prop:pet-ch7-busemann-approximants}
  \dcref{ch7:7.3.9}
  If $\Ric(M,g)\ge0$, then $\Delta b_c\le0$ everywhere.
\end{lemma}
\begin{proof}
  \uses{lem:pet-ch7-busemann-superharmonic}
  \sketch
  By \cref{prop:pet-ch7-busemann-asymptote-relation}, \(b_c(p)+b_{\widetilde c}\) is an upper support for \(b_c\) at \(p\). The approximants \(b_t=|\cdot\,\widetilde c(t)|-t\) are smooth upper supports with \(\Delta b_t(p)\le(n-1)/t\). Letting \(t\to\infty\) gives \(\Delta b_c(p)\le0\).
\end{proof}

\begin{theorem}[Cheeger--Gromoll splitting theorem]
  \label{thm:pet-ch7-splitting-theorem}
  \lean{PetersenLib.cheegerGromollSplittingTheorem}
  \source{petersen:page-0315,petersen:page-0318,petersen:page-0319}
  \uses{lem:pet-ch7-busemann-superharmonic, thm:pet-ch7-strong-maximum-principle, thm:pet-ch7-harmonic-regularity}
  \dcref{ch7:7.3.5}
  If $(M,g)$ contains a line and has $\Ric\ge0$, then $(M,g)$ is isometric to
  a product $(H\times\mathbb{R},\,g_0+dt^2)$.
\end{theorem}
\begin{proof}
  \uses{thm:pet-ch7-splitting-theorem,lem:pet-ch7-busemann-superharmonic,thm:pet-ch7-strong-maximum-principle,thm:pet-ch7-harmonic-regularity,prop:pet-ch7-busemann-approximants,def:pet-ch7-busemann-function,prop:pet-ch7-volume-form-laplacian-trace}
  \sketch
  For a line \(c\), \cref{lem:pet-ch7-busemann-superharmonic} makes \(b^++b^-\) superharmonic; it is nonnegative and vanishes on \(c\), so the maximum principle makes it zero. Thus each function is weakly harmonic. Support inequalities give \(|\nabla b^\pm|=1\), and \cref{thm:pet-ch7-harmonic-regularity} makes them smooth. The trace identity with \(\Delta b^+=0\) forces \(\Hess b^+=0\), yielding \(H\times\mathbb R\).
\end{proof}

\subsection{Structure Results in Nonnegative Ricci Curvature}

The splitting theorem gives several structural consequences for compact
manifolds of nonnegative Ricci curvature.

\begin{corollary}[Ricci-flat obstruction for sphere products]
  \label{cor:pet-ch7-no-ricci-flat-sphere-product}
  \lean{PetersenLib.noRicciFlatMetricOnSphereProductLowDim}
  \source{petersen:page-0319}
  \uses{thm:pet-ch7-splitting-theorem}
  \dcref{ch7:7.3.10}
  $S^k\times S^1$ admits no Ricci-flat metric when $k=2,3$.
\end{corollary}
\begin{proof}
  \uses{cor:pet-ch7-no-ricci-flat-sphere-product}
  \sketch
  The universal cover is \(S^k\times\mathbb R\), which is disconnected at infinity; the splitting theorem therefore gives a Ricci-flat compact simply connected \(k\)-manifold homotopy equivalent to \(S^k\). For \(k=2,3\) it is flat, contradicting the fundamental group of a compact flat manifold.
\end{proof}

\begin{theorem}[structure of compact manifolds with nonnegative Ricci curvature]
  \label{thm:pet-ch7-structure-theorem}
  \lean{PetersenLib.nonnegativeRicciStructureTheorem}
  \source{petersen:page-0319,petersen:page-0320}
  \uses{thm:pet-ch7-splitting-theorem, lem:pet-ch7-ray-line-existence}
  \dcref{ch7:7.3.11}
  Suppose $(M,g)$ is compact with $\Ric\ge0$. Then:
  \begin{enumerate}
    \item[(1)] the universal cover $(\tilde M,\tilde g)$ splits isometrically
      as $N\times\mathbb{R}^k$, with $N$ compact and containing no lines;
    \item[(2)] $\Iso(\tilde M)=\Iso(N)\times\Iso(\mathbb{R}^k)$;
    \item[(3)] there is a finite normal subgroup $G\subset\pi_1(M)$ with
      $\pi_1(M)/G\cong\pi_1(M)\cap\Iso(\mathbb{R}^k)$, and a finite-index subgroup
      $\mathbb{Z}^k\subset\pi_1(M)\cap\Iso(\mathbb{R}^k)$.
  \end{enumerate}
\end{theorem}
\begin{proof}
  \uses{thm:pet-ch7-structure-theorem}
  \sketch
  Iterate the splitting theorem on the universal cover to obtain \(N\times\mathbb R^k\) with \(N\) line-free; compactness of \(N\) follows by recentering a ray under the cocompact deck action, which would otherwise produce a line. Isometries preserve the Euclidean line factors, so the isometry group splits. The kernel of the Euclidean projection is finite, and the translation subgroup is a discrete cocompact lattice, hence a finite-index \(\mathbb Z^k\).
\end{proof}

\begin{remark}[Wilking's converse realization theorem]
  \label{rem:pet-ch7-wilking-remark}
  \lean{PetersenLib.wilkingConverseRemark}
  \source{petersen:page-0320}
  \uses{thm:pet-ch7-structure-theorem}
  \dcref{ch7:7.3.12}
  Wilking has shown conversely that any group $G$ admitting a finite normal
  subgroup $H\subset G$ with $G/H$ acting discretely and cocompactly on a
  Euclidean space is the fundamental group of a compact manifold with
  nonnegative sectional curvature.
\end{remark}

\begin{corollary}[aspherical nonnegative-Ricci manifolds are flat]
  \label{cor:pet-ch7-flat-manifold-Kpi1}
  \lean{PetersenLib.flatIfAsphericalNonnegRicci}
  \source{petersen:page-0320,petersen:page-0321}
  \uses{thm:pet-ch7-structure-theorem}
  \dcref{ch7:7.3.13}
  Suppose $(M,g)$ is compact with $\Ric\ge0$. If $M$ is a $K(\pi,1)$ (its
  universal cover is contractible), then the universal cover is Euclidean
  space and $(M,g)$ is flat.
\end{corollary}
\begin{proof}
  \uses{cor:pet-ch7-flat-manifold-Kpi1}
  \sketch
  By \cref{thm:pet-ch7-structure-theorem}, \(\widetilde M=N\times\mathbb R^k\) with \(N\) compact. Contractibility of \(\widetilde M\) forces \(N\) to be a point, so \(\widetilde M=\mathbb R^n\) and the metric is flat.
\end{proof}

\begin{corollary}[positive Ricci curvature at one point forces finite fundamental group]
  \label{cor:pet-ch7-finite-fundamental-group-positive-ricci}
  \lean{PetersenLib.finitePi1IfPositiveRicciSomewhere}
  \source{petersen:page-0321}
  \uses{thm:pet-ch7-structure-theorem}
  \dcref{ch7:7.3.14}
  If $(M,g)$ is compact with $\Ric\ge0$ and $\Ric>0$ on some tangent space
  $T_pM$, then $\pi_1(M)$ is finite.
\end{corollary}
\begin{proof}
  \uses{cor:pet-ch7-finite-fundamental-group-positive-ricci}
  \sketch
  A positive Ricci tensor at one point is incompatible with a nontrivial Euclidean factor in the structure decomposition, since that direction has zero Ricci curvature. Thus the universal cover is compact and the deck group is finite.
\end{proof}

\begin{corollary}[first Betti-number bound]
  \label{cor:pet-ch7-first-betti-number-bound}
  \lean{PetersenLib.firstBettiNumberBoundNonnegRicci}
  \source{petersen:page-0321}
  \uses{thm:pet-ch7-structure-theorem}
  \dcref{ch7:7.3.15}
  If $(M,g)$ is compact with $\Ric\ge0$, then $b_1(M)\le\dim M=n$, with
  equality iff $(M,g)$ is a flat torus.
\end{corollary}
\begin{proof}
  \uses{cor:pet-ch7-first-betti-number-bound}
  \sketch
  The structure theorem gives a finite quotient whose Euclidean subgroup contains a finite-index \(\mathbb Z^k\). Abelianization therefore surjects onto \(\mathbb Z^{b_1}\), giving \(b_1\le k\le n\). If \(b_1=n\), then \(k=n\), the compact factor and finite kernel vanish, and the deck group is a lattice; hence \(M\) is a flat torus.
\end{proof}

\begin{theorem}[homogeneous nonnegative-Ricci splitting]
  \label{thm:pet-ch7-homogeneous-splitting}
  \lean{PetersenLib.homogeneousNonnegRicciSplitting}
  \source{petersen:page-0321,petersen:page-0322}
  \uses{thm:pet-ch7-splitting-theorem, thm:pet-ch7-structure-theorem}
  \dcref{ch7:7.3.16}
  Let $(M,g)$ be a homogeneous Riemannian manifold with $\Ric\ge0$. Then
  $(M,g)=(N\times\mathbb{R}^k,\,g_N+g_{\mathbb{R}^k})$ with $(N,g_N)$ compact homogeneous.
\end{theorem}
\begin{proof}
  \uses{thm:pet-ch7-homogeneous-splitting}
  \sketch
  Apply the structure theorem to split \(M=N\times\mathbb R^k\) with \(N\) line-free. Homogeneity preserves this product and makes \(N\) homogeneous. A ray in a noncompact homogeneous \(N\), recentered by isometries, has a subsequential backward extension to a line, contradicting the choice of \(N\); therefore \(N\) is compact.
\end{proof}

\section{Exercises}

\begin{remark}[polar-density comparison equation]
  \label{rem:pet-ch7-ex-1}
  \lean{PetersenLib.exercise7_5_1}
  \source{petersen:page-0322}
  \uses{prop:pet-ch7-volume-form-laplacian-trace, def:pet-ch7-polar-density}
  \dcref{ch7:7.5.1}
  With notation as in §7.1.1 and using $\vol=\lambda\,dr\wedge\vol_{n-1}$, show
  that $\mu=\lambda^{1/(n-1)}$ satisfies
  \[
    \partial_r^2\mu\le-\frac{\mu}{n-1}\Ric(\partial_r,\partial_r),\qquad
    \mu(0,\theta)=0,\qquad \lim_{r\to0}\partial_r\mu(r,\theta)=1,
  \]
  and that this yields the desired volume-form estimates.
\end{remark}

\begin{remark}[Calabi--Yau linear volume growth]
  \label{rem:pet-ch7-ex-2}
  \lean{PetersenLib.exercise7_5_2}
  \source{petersen:page-0322}
  \uses{thm:pet-ch7-structure-theorem}
  \dcref{ch7:7.5.2}
  Let $(M,g)$ be complete noncompact with $\Ric\ge0$, $p\in M$.
  \begin{enumerate}
    \item[(1)] Show that for each $R>1$ there is $x\in M$ with
    \[
      \vol B(p,1)\le\vol B(x,R+1)-\vol B(x,R-1)
      \le\frac{(R+1)^n-(R-1)^n}{(R+1)^n}\vol B(p,2R).
    \]
    \item[(2)] Show there is $C>0$ with $\vol B(p,R)\ge CR$.
  \end{enumerate}
\end{remark}

\begin{remark}[equivalent definitions of convexity on an interval]
  \label{rem:pet-ch7-ex-3}
  \lean{PetersenLib.exercise7_5_3}
  \source{petersen:page-0322,petersen:page-0323}
  \uses{def:pet-ch7-support-function}
  \dcref{ch7:7.5.3}
  Let $f:I\to\mathbb{R}$ be continuous on an interval $I\subset\mathbb{R}$. Show that (1)
  $f$ is convex, (2) $f$ has a linear support function from below of the form
  $a(x-x_0)+f(x_0)$ at every $x_0\in I$, and (3) $f''\ge0$ in the support
  sense at every $x_0\in I$, are equivalent.
\end{remark}

\begin{remark}[failure of an endpoint Poincaré estimate]
  \label{rem:pet-ch7-ex-4}
  \lean{PetersenLib.exercise7_5_4}
  \source{petersen:page-0323}
  \uses{thm:pet-ch7-poincare-sobolev}
  \dcref{ch7:7.5.4}
  Show that on a compact Riemannian manifold there is no finite $S(s)$ with
  $\|f-f_M\|_\infty\le S\|df\|_1$ for $1<s<\dim M$.
\end{remark}

\begin{remark}[basic covering lemma]
  \label{rem:pet-ch7-ex-5}
  \lean{PetersenLib.exercise7_5_5}
  \source{petersen:page-0323}
  \uses{thm:pet-ch7-maximal-function-theorem}
  \dcref{ch7:7.5.5}
  Given a separable metric space $(X,d)$ and a bounded positive function
  $R:X\to(0,D]$, show there is a countable $A\subset X$ such that the balls
  $B(p,R(p))$, $p\in A$, are pairwise disjoint and $X=\bigcup_{p\in
  A}B(p,5R(p))$. (Hint: choose the points successively so that
  $R(p_{k+1})\ge\tfrac12\sup_{p\in X\setminus\bigcup_iB(p_i,2R(p_i))}R(p)$.)
\end{remark}

\begin{remark}[radial Hessian rigidity]
  \label{rem:pet-ch7-ex-6}
  \lean{PetersenLib.exercise7_5_6}
  \source{petersen:page-0323}
  \uses{thm:pet-ch7-cheng-rigidity}
  \dcref{ch7:7.5.6}
  Assume $r(x)=|xp|$ is smooth on $B(p,R)$. Show that if
  $\Hess r=(\snk'(r)/\snk(r))g_r$ in polar coordinates, then all sectional
  curvatures on $B(p,R)$ equal $k$.
\end{remark}

\begin{remark}[unbounded Sobolev--Poincaré constants]
  \label{rem:pet-ch7-ex-7}
  \lean{PetersenLib.exercise7_5_7}
  \source{petersen:page-0323}
  \uses{thm:pet-ch7-poincare-sobolev}
  \dcref{ch7:7.5.7}
  Construct convex surfaces in $\mathbb{R}^3$ by capping off cylinders
  $[-R,R]\times S^1$ to show the Sobolev--Poincaré constants increase as $R$
  increases. (Hint: consider test functions constant except on $[-1,1]\times
  S^1$.)
\end{remark}

\begin{remark}[rigidity in absolute volume comparison]
  \label{rem:pet-ch7-ex-8}
  \lean{PetersenLib.exercise7_5_8}
  \source{petersen:page-0323}
  \uses{lem:pet-ch7-absolute-volume-comparison}
  \dcref{ch7:7.5.8}
  Show that if $(M,g)$ has $\Ric\ge(n-1)k$ and $\vol B(p,R)=v(n,k,R)$ for some
  $p\in M$, then the metric has constant curvature $k$ on $B(p,R)$.
\end{remark}

\begin{remark}[Jacobi fields and the Jacobian of an exponential flow]
  \label{rem:pet-ch7-ex-9}
  \lean{PetersenLib.exercise7_5_9}
  \source{petersen:page-0323}
  \uses{lem:pet-ch7-laplacian-comparison-calabi}
  \dcref{ch7:7.5.9}
  Let $X$ be a vector field and $F_t(p)=\exp_p(tX|_p)$.
  \begin{enumerate}
    \item[(1)] For $v\in T_pM$, show $J(t)=DF_t(v)$ is a Jacobi field along
      $c(t)=\exp(tX)$ with $J(0)=v$, $\dot J(0)=\nabla_vX$.
    \item[(2)] For an orthonormal basis $e_i$ of $T_pM$ and $J_i(t)=DF_t(e_i)$,
      show $(\det[DF_t])^2=\det[g(J_i(t),J_j(t))]$.
    \item[(3)] Show that as long as $\det(DF_t)\ne0$,
      \[
        \frac{d^2(\det(DF_t))^{1/n}}{dt^2}
        \le-\frac{(\det(DF_t))^{1/n}}{n}\Ric(\dot c,\dot c).
      \]
      (Hint: use $(\operatorname{tr}A)^2\le n\operatorname{tr}(A^*A)$ for
      $n\times n$ matrices $A$.)
  \end{enumerate}
\end{remark}

\begin{remark}[Euclidean rigidity from asymptotic volume ratio]
  \label{rem:pet-ch7-ex-10}
  \lean{PetersenLib.exercise7_5_10}
  \source{petersen:page-0324}
  \uses{lem:pet-ch7-relative-volume-comparison}
  \dcref{ch7:7.5.10}
  Show that a complete manifold $(M,g)$ with $\Ric\ge0$ and
  $\lim_{r\to\infty}\vol B(p,r)/(\omega_nr^n)=1$ for some $p\in M$ must be
  isometric to Euclidean space.
\end{remark}

\begin{remark}[trace inequality for the Hessian]
  \label{rem:pet-ch7-ex-11}
  \lean{PetersenLib.exercise7_5_11}
  \source{petersen:page-0324}
  \uses{lem:pet-ch7-hessian-touching-comparison}
  \dcref{ch7:7.5.11}
  Show that any function on an $n$-dimensional Riemannian manifold satisfies
  $|\Hess u|^2\ge\tfrac1n|\Delta u|^2$, with equality iff
  $\Hess u=\tfrac{\Delta u}{n}g$. What can be said about $M$ when
  $\Hess u=\tfrac{\Delta u}{n}g$?
\end{remark}

\begin{remark}[integration by parts for the Laplacian]
  \label{rem:pet-ch7-ex-12}
  \lean{PetersenLib.exercise7_5_12}
  \source{petersen:page-0324}
  \uses{def:pet-ch7-linear-harmonic-function}
  \dcref{ch7:7.5.12}
  Show that if $u,v:M\to\mathbb{R}$ are compactly supported and smooth on open dense
  sets, then
  \[
    \int u\,\Delta v\,\vol=\int v\,\Delta u\,\vol
    =-\int g(du,dv)\,\vol=-\int g(\nabla u,\nabla v)\,\vol.
  \]
\end{remark}

\begin{remark}[sign of Laplacian eigenvalues]
  \label{rem:pet-ch7-ex-13}
  \lean{PetersenLib.exercise7_5_13}
  \source{petersen:page-0324}
  \uses{thm:pet-ch7-strong-maximum-principle}
  \dcref{ch7:7.5.13}
  Show that if $\Delta u=\lambda u$ on a closed Riemannian manifold, then
  $\lambda\le0$, and $\lambda=0$ forces $u$ constant.
\end{remark}

\begin{remark}[a first spherical eigenfunction]
  \label{rem:pet-ch7-ex-14}
  \lean{PetersenLib.exercise7_5_14}
  \source{petersen:page-0324}
  \uses{thm:pet-ch7-cheng-rigidity}
  \dcref{ch7:7.5.14}
  Show that $u_k=\cos(\sqrt kr)$ on $S^n_k=S^n(1/\sqrt k)$ satisfies $\Delta
  u_k=-nk\,u_k$ and $\int u_k\,\vol=0$.
\end{remark}

\begin{remark}[Lichnerowicz eigenvalue estimate]
  \label{rem:pet-ch7-ex-15}
  \lean{PetersenLib.exercise7_5_15}
  \source{petersen:page-0324}
  \uses{thm:pet-ch7-cheng-rigidity}
  \dcref{ch7:7.5.15}
  Let $(M^n,g)$ be closed with $\Ric\ge(n-1)k>0$. Use the Bochner formula to
  show that every $u$ with $\Delta u=-\lambda u$, $\lambda>0$, satisfies
  $\lambda\ge nk$. Deduce, via the spectral theorem for $\Delta$, that every
  $u$ with $\int u\,\vol=0$ satisfies the Poincaré inequality
  $\int u^2\,\vol\le\tfrac1{nk}\int|du|^2\,\vol$.
\end{remark}

\begin{remark}[Obata rigidity]
  \label{rem:pet-ch7-ex-16}
  \lean{PetersenLib.exercise7_5_16}
  \source{petersen:page-0324}
  \uses{rem:pet-ch7-ex-15, thm:pet-ch7-cheng-rigidity}
  \dcref{ch7:7.5.16}
  Let $(M^n,g)$ be closed with $\Ric\ge(n-1)k>0$. Using the Bochner formula as
  in \cref{rem:pet-ch7-ex-15}, show that if $\Delta u=-nk\,u$ for some
  function $u$, then $\Hess u=-ku\,g$. Conclude $(M^n,g)=S^n_k$.
\end{remark}

\begin{remark}[Li--Schoen local Poincaré estimate]
  \label{rem:pet-ch7-ex-17}
  \lean{PetersenLib.exercise7_5_17}
  \source{petersen:page-0324,petersen:page-0325}
  \uses{thm:pet-ch7-poincare-sobolev}
  \dcref{ch7:7.5.17}
  Let $(M,g)$ be complete with $\Ric\ge-(n-1)k^2$, $k\ge0$; $p\in M$; $R>0$
  with $\partial B(p,2R)\ne\emptyset$; $q\in\partial B(p,2R)$;
  $r(x)=|xq|$.
  \begin{enumerate}
    \item[(1)] Show $\Delta r\le(n-1)(R^{-1}+k)$ on $B(p,R)$.
    \item[(2)] Let $f(x)=a\exp(-ar(x))$, $a>0$. Show
      $\Delta f\ge a\exp(-3aR)\big(a-(n-1)(R^{-1}+k)\big)$.
    \item[(3)] Let $u\ge0$ be smooth with compact support in $B(p,R)$, and
      choose $a=n(R^{-1}+k)$. Using
      $\int_{B(p,R)}u\,\Delta f\,\vol=-\int_{B(p,R)}g(du,df)\,\vol$, show
      $\int_{B(p,R)}u\,\vol\le C\int_{B(p,R)}|du|\,\vol$, where
      $C=\dfrac{R}{1+kR}\exp\big(2n(1+kR)\big)$.
    \item[(4)] Extend the inequality of (3) to all smooth $u$ with compact
      support in $B(p,R)$.
    \item[(5)] For $s\ge1$ and $u$ with compact support in $B(p,R)$, show
      $\int_{B(p,R)}|u|^s\,\vol\le(sC)^s\int_{B(p,R)}|du|\,\vol$.
  \end{enumerate}
\end{remark}

\begin{remark}[Cheeger comparison for unions of balls]
  \label{rem:pet-ch7-ex-18}
  \lean{PetersenLib.exercise7_5_18}
  \source{petersen:page-0325}
  \uses{lem:pet-ch7-relative-volume-comparison}
  \dcref{ch7:7.5.18}
  Suppose $(M^n,g)$ has $\Ric\ge(n-1)k$.
  \begin{enumerate}
    \item[(1)] For $p_1,\dots,p_k\in M$, show
      $r\mapsto\vol\big(\bigcup_{i=1}^kB(p_i,r)\big)/v(n,k,r)$ is
      nonincreasing and converges to $k$ as $r\to0$.
    \item[(2)] For $A\subset M$, show
      $r\mapsto\vol\big(\bigcup_{p\in A}B(p,r)\big)/v(n,k,r)$ is
      nonincreasing, by applying (1) to a dense finite subset of $A$.
  \end{enumerate}
\end{remark}

\begin{remark}[polar-cone volume comparison]
  \label{rem:pet-ch7-ex-19}
  \lean{PetersenLib.exercise7_5_19}
  \source{petersen:page-0326}
  \uses{lem:pet-ch7-absolute-volume-comparison}
  \dcref{ch7:7.5.19}
  For $p\in M$ and $\Gamma\subset T_pM$ a set of unit vectors, define the
  polar cone $B^\Gamma(p,R)=\{(t,\theta)\in M\mid t\le R,\ \theta\in\Gamma\}$.
  If $\Ric M\ge(n-1)k$, show
  \[
    \vol B^\Gamma(p,R)\le\vol\Gamma\cdot\int_0^R(\snk_k(t))^{n-1}\,dt.
  \]
\end{remark}

\begin{remark}[fundamental group and center of a compact Lie group]
  \label{rem:pet-ch7-ex-20}
  \lean{PetersenLib.exercise7_5_20}
  \source{petersen:page-0326}
  \uses{thm:pet-ch7-structure-theorem}
  \dcref{ch7:7.5.20}
  Let $G$ be a compact connected Lie group with a biinvariant metric and
  $\Ric\ge0$.
  \begin{enumerate}
    \item[(1)] If $G$ has finite center, show $\pi_1(G)$ is finite.
    \item[(2)] Show a finite cover of $G$ has the form $G'\times T^k$ with
      $G'$ compact simply connected and $T^k$ a torus.
    \item[(3)] If $\pi_1(G)$ is finite, show the center of $G$ is finite.
  \end{enumerate}
\end{remark}

\begin{remark}[Euclidean-end rigidity]
  \label{rem:pet-ch7-ex-21}
  \lean{PetersenLib.exercise7_5_21}
  \source{petersen:page-0326}
  \uses{thm:pet-ch7-structure-theorem}
  \dcref{ch7:7.5.21}
  Let $(M,g)$ be $n$-dimensional and isometric to Euclidean space outside a
  compact $K\subset M$, i.e.\ $M-K$ is isometric to $\mathbb{R}^n-C$ for some
  compact $C\subset\mathbb{R}^n$. If $\Ric\ge0$, show $M=\mathbb{R}^n$.
\end{remark}

\begin{remark}[Myers diameter bound from excess]
  \label{rem:pet-ch7-ex-22}
  \lean{PetersenLib.exercise7_5_22}
  \source{petersen:page-0326}
  \uses{thm:pet-ch7-cheng-rigidity}
  \dcref{ch7:7.5.22}
  Show that if $\Ric\ge n-1$ then $\diam\le\pi$, by showing that if
  $|pq|>\pi$, then $e_{p,q}(x)=|px|+|xq|-|pq|$ has negative Laplacian at a
  local minimum.
\end{remark}
